\documentclass[a4paper,11pt]{article}

\usepackage{revy}
\usepackage[utf8]{inputenc}
\usepackage[T1]{fontenc}
\usepackage[danish]{babel}


\revyname{MatematikRevyen}
\revyyear{2026}
% HUSK AT OPDATERE VERSIONSNUMMER
\version{0.1}
\eta{$6$ minutter}
\status{Ikke færdig}

\title{Skal vi køre stærkt?}
\author{Pernille '25 \& Niklas '21}

\begin{document}
\maketitle

\begin{roles}
\role{R}[Skuespiller] "Rus" 
\role{S1}[Skuespiller] "Ældre studerende 1"
\role{S2}[Skuespiller] "Ældre studerende 2"
\role{K1}[Statist] "Guld kapsel"
\role{K2}[Statist] "Grøn kapsel"
\role{K3}[Statist] "Cola kapsel"
\role{F}[Statist] "Fysiker/politbetjent"
\end{roles}

\begin{props}
\prop{Racerbilebane med et fængselsfelt på felt 12}
\prop{3 kapselhjelme, en guld, en grøn og en cola}
\prop{Håndjern}
\prop{Grene}
\prop{3 terninger}
\prop{Fez}
\prop{grøn body suite}
\prop{3 øl, S1, S2 og R drikker af}
\prop{Cafeen? baggrund}
\end{props}



\begin{sketch}

\scene{S1 og S2 sidder omkring et bord på cafeen?. Der er god og festlig stemning en fredag aften med musik i baggrunden. Der er en kæmpe racerbilebane rundt om dem. R kommer hen til bordet.}

\says{R} Hvad spiller I?

\says{S1} Vi skal lige til at spille racerbile. Skal du være med?

\says{R} Racebiler? Hvordan spiller man det? 

\says{S2} Det er ikke så svært. Du lærer reglerne undervejs

\says{R} Okay. Lad os!

\scene{R sætter sig ved bordet.}

\says{S1} Vi skal hver især vælge en kapsel som brik. Jeg tager guldkapslen. 

\scene{K1 iført et Guld kapsel kostume kommer ind og stiller sig på banen}

\says{S2} Så vælger jeg grøn.

\scene{K2 iført et Tubrog Grøn kapsel kostume kommer ind og stiller sig på banen.}

\says{R} Så er der kun cola kaplsen tilbage, så den er jeg.

\scene{K3 iført et Cola kapsel kostume kommer ind og stiller sig på banen.}

\says{S1} Jeg starter. 

\scene{S1 slår med terningerne}

\says{S1} En 6'er. 

\scene{S1 kigger lige mod publikum i et par sekunder (Det må man ikke i det rigtige racerbile)}

\says{S1} Ja, det må jeg gerne. 

\scene{K1 rykker 6 felter frem imens spillerne kigger på den.}

\says{S1} Og så skal jeg drikke en straftår, fordi jeg er den ældste mat'er her. 

\says{S2} Min tur.

\scene{S2 tager terningerne og slår.}

\says{S1} En 1'er

\says{R} Øv hva? Så kommer du ikke så langt på banen? 

\says{S2} Ha! Dumme rus. For JEG har den her.

\scene{S2 Trækker et kæmpe plus 4 kort ud af trøjen og K2 rykker 5 frem.}

\says{S2} Og jeg skal selvfølgelig ikke drikke nogle straftårer, men rus drikker 4 plus. 

\scene{R drikker uforstående sine tårer og tager terningerne.}

\says{S1} Hov ikke så hurtigt. Du må da ikke bare tage terningerne! 

\says{R} Men det måtte S2 da? 

\says{S2} Det er fordi, jeg bestod sand i 1. forsøg. 

\says{R} Nååårh. Jaa. Okay så. 

\scene{R tager en tår og slår med terningen.}

\says{R} En 3'er. 

\scene{ S1 og S2 rejser sig hastigt op og stiller sig med armene ud som et træ. (evt. har de grene i hånden.)}

\says{S1 og S2} [Råber] TRÆ!

\scene{Både R og K3 bliver forskrækket. S1 og S2 sætter sig hurtigt ned igen og peger på R}

\says{S2} [peger på R] Drik!

\scene{R drikker forvirret sin tår. K3 begynder forsigtigt at gå 3 felter fremad, imens S1 og S2 holder grundigt øje.}

\says{R} Skal jeg drikke nogle straftårer for noget nu? 

\scene{S1 og S2 lades som om, de ikke ved det. }

\says{R} Så er min tur slut.

\says{S1 og S2}[Råber i kor] FEJL! 

\says{S1} Du skal drikke for at slå skarpt mindre end 4. 

\says{S2} Og for ikke at begynde med højre ben og slutte med venstre. 

\scene{De tænker begge to for at opfinde en grund mere.}

\says{S1} Og... og...  fordi du er rus. 

\says{S2} Ja! Så drikker du dobbelt.

\scene{Rus er forvirret men drikker.}

\scene{S1 tager terningerne og slår mnu ed 2 terninger.}

\says{S1} En 8'er. Det er katagorirunde. Jeg vælger katagorien KomAn 2 sætninger!

\says{R} Hvem kommer an? Imod os?  og hvorfor 2 gange? 

\says{S1} Riemann mapping theorem.

\says{S2} Montel's sætning.

\says{R}[Frustreret og forvirret.] Hvad? 

\says{S1} "Hvad" theoremet? Den kender jeg ikke.

\says{S2} Du skal drikke!

\scene{K1 rykker 8 felter.}

\says{S1} Og jeg skal ikke drikke straftårer for noget.

\says{R} Ha! du skal jo drikke for at være ældst!

\says{S1} Nårh ja! Godt du siger det. Det skifter til yngst i denne runde, så du skal drikke!

\scene{Rus drikker opgivende og frustreret igen.}

\scene{S2 tager terningerne og slår med 2 terninger.}

\says{S2} En 13'er! Så skal alle vise Greens sætning.

\scene{S1 og S2 kigger på hinanden kort.}

\says{S1} Altså jeg viste den for mange år siden.

\says{S2} Det gjorde jeg også.

\scene{S1 og S2 kigger på R}

\says{R} Øhhh hvad? Hvad er det? Det kan jeg da ikke.

\says{S1} Så skal du have en lille fysiker. 

\scene{En statist iført grønt bodysuit og Fez kommer ind fra sidetæppet.}

\says{R} Ej det er simpelthen for meget. Det kan da ikke være rigtigt. Jeg skal fandme ikke have nogen fysiker. 

\scene{Statisten går skuffet ud igen og  K2 rykker 13 frem.}

\says{S2} Okay okay. Du slipper denne gang, men det er din tur. 

\scene{Lidt modvilligt slår R med 2 terninger} 

\says{R} 9. 

\scene{K3 rykker frem og rammer fængselsfeltet.}

\says{S1} ååhhh!!! Du kommer i fængsel!

\scene{Politibetjent kommer ind, eventuelt grøn statist med politihat på i stedet for fez og jakke. Han lægger K3 ned og lægger K3 i håndjern imens K3 kæmper imod.}

\says{R} Hvad fanden sker der? Lad mig gætte. Jeg skal bunde, imens jeg sidder i en børnepool i speedos eller sådan noget dumt? 

\says{S2} Nej nej fjolle. Det er kun, hvis du havde slået 15. Vi skal skåle for dig i solidaritet, imens du afsoner din dom. 

\scene{S1 og S2 Skåler og tager en tår. S1 tager nu terningerne og slår 12.}  

\says{S1} 12

\scene{S1 og S2 kigger begge samtidig over på R og venter på R gør noget}

\says{R} Hvad sker der?

\says{S2} S1 slog 12, og der er fuldmåne, så personen til venstre skal hyle som en ulv. Det gjorde du ikke, så nu skal du bunde. 

\says{R}[Rejser sig op] Ej nu kan det være nok. Det giver jo ingen mening det her. Jeg gider ikke være med mere. 

\says{S1} Rolig nu. Vi er jo så godt i gang

\says{R} Ej jeg går nu. 

\scene{R rejser sig op.}

\says{S1} Okay okay. Du behøver ikke gå Så lad os spille noget andet.

\says{R} Okay så. Hvad tænker I?

\says{S2} Hvad med Kranen?

\says{R} Kranen? Hvordan spiller man det? 

\scene{Lys ned.}

\end{sketch}
\end{document}


Mere skarp slutning


alternativ slutning. Rus går og så snakker de studerende om at de faktisk heller ikke kender reglerne og de bare fandt på det undervejs.